\section{Software Analysis and Program Comprehension}

Software analysis is a broad term that covers techniques used to inspect,
summarize, evaluate, and reason about software artifacts. In the context of this
project, the term includes static analysis, visual analysis, machine learning
over source artifacts, and tool-supported feedback inside development
environments. These activities are closely related to program comprehension:
before a developer can modify, review, test, or refactor a system, the developer
must form a working understanding of what its artifacts are, how they are
organized, and what roles they play.

Program comprehension is difficult because source code is not only text. It is
also structure, convention, history, style, architecture, dependency, and domain
knowledge. The same syntactic construct may have different meaning depending on
frameworks and project conventions; conversely, visually similar files may
perform similar roles even when they belong to different projects. This tension
motivates the search for representations that capture useful structural signals
without requiring a full semantic interpretation at the first step.

Traditional program-comprehension support includes code navigation,
cross-reference search, syntax highlighting, dependency views, metrics,
architectural diagrams, and static-analysis warnings. These mechanisms are
essential, but they tend to operate either at the textual level or through
language-aware abstractions. A minimap-based approach adds another layer: it
represents a source artifact as an image-like signature that can be inspected by
models and, in some cases, by humans. It does not aim to answer every semantic
question about a file. Instead, it supports a preliminary form of recognition:
what type of artifact does this resemble, which project conventions does it
share, and which visual regions explain a classification?

\section{Static Analysis and Source Code Representations}

Static analysis examines software artifacts without executing them. It may
detect coding-rule violations, compute metrics, identify possible defects,
extract dependencies, or enforce architectural constraints. Static analysis is
particularly valuable because it can run automatically and repeatedly during
development, testing, code review, and continuous integration. However, static
analysis usually requires parsers, language grammars, rule definitions, or
framework-specific knowledge.

Source code can be represented in several ways for static or learning-based
analysis. Table~\ref{tab:foundations-representations} summarizes common
representations and their implications for this project.

\begin{table}[ht]
\centering
\caption{Common source code representations and their implications.}
\label{tab:foundations-representations}
\begin{tabular}{p{0.21\textwidth}p{0.35\textwidth}p{0.31\textwidth}}
\toprule
\textbf{Representation} & \textbf{Strengths} & \textbf{Limitations} \\
\midrule
Raw text & Simple, universal, preserves full content. & Sensitive, noisy, and difficult to compare across languages. \\
Tokens & Captures lexical structure and supports sequence models. & Requires tokenization and often preserves identifiers. \\
Abstract syntax trees & Captures grammar and hierarchy. & Requires language-specific parsers and valid code. \\
Graphs & Represents dependencies, flows, and structural relations. & Can be expensive to build and normalize. \\
Metrics & Compact and interpretable. & Often loses local structure and layout. \\
Images/minimaps & Compact, visual, language-light, compatible with computer vision. & Loses explicit semantics and may require careful interpretation. \\
\bottomrule
\end{tabular}
\fautor
\end{table}

Learning-based software engineering has explored several of these
representations. Surveys and empirical studies discuss how tokens, trees,
graphs, and neural representations can support code classification,
summarization, defect prediction, and related tasks
\cite{allamanis2018survey,alon2019code2vec,feng2020codebert,guo2021graphcodebert}.
Graph-based representations are especially useful when dependencies and program
structure are central, as in just-in-time bug prediction using source code graphs
\cite{nadim2022sourcecodegraphs}. Sparse and transformer-based models seek to
handle long code sequences more efficiently \cite{yang2025sparsecoder}.

The minimap representation proposed in this project is not intended to replace
these approaches. Instead, it provides a complementary abstraction. A minimap is
less semantically explicit than an abstract syntax tree or graph, but it is
easier to generate for heterogeneous file types, including configuration files,
markup, scripts, and generated artifacts. It can also be anonymized in a way that
removes identifiers and literal values while retaining layout and structural
symbols.

\section{Software Visualization}

Software visualization studies techniques for representing software artifacts,
relationships, evolution, and runtime or static properties in visual forms.
Visualization can support program comprehension by reducing cognitive load,
revealing patterns, and allowing analysts to reason about large systems at
different levels of abstraction.

Several software visualization tools focus on structural understanding.
Pinzger et al. \cite{pinzger2008tool} presented a tool for visual understanding
of source code dependencies, emphasizing the role of visual relationships in
program comprehension. SourceMiner Evolution supports the comprehension of
feature evolution through coordinated views such as treemaps, polymetric views,
and dependency graphs \cite{novais2013sourceminer}. SArF Map visualizes software
architecture from feature and layer viewpoints, using map metaphors to help
analysts inspect implicit structures \cite{kobayashi2013sarf}.

More recent work extends this landscape. FineCodeAnalyzer provides
multi-perspective source code analysis through fine-grained interactive
visualization \cite{qayum2022finecodeanalyzer}. SFLVis supports visual analysis
for software fault localization \cite{sun2024sflvis}. Code-smell detection and
visualization have also been systematized in the literature, showing that visual
support remains relevant for understanding quality problems \cite{reis2022code}.
Tools such as CodePanorama further explore zoomed-out and language-agnostic code
inspection \cite{etter2022codepanorama}.

This project relates to software visualization but shifts the main audience of
the visual artifact. Traditional visualizations are usually designed for human
interpretation. Source code minimaps, in this project, are designed primarily as
machine-readable visual signatures. Humans may inspect them, and explainability
methods may expose salient regions, but the core pipeline treats minimaps as
inputs for computational models.

\section{Source Code as Visual Data}

Representing non-visual phenomena as images is not new in machine learning.
Audio signals, for example, can be transformed into spectrograms and then
processed with image-analysis techniques. Costa et al. used spectrograms for
music genre recognition \cite{costa2011music}, and later work showed that Local
Binary Patterns can be effective texture descriptors for music classification
\cite{costa2012music}. Convolutional neural networks have also been evaluated
for music classification using spectrograms \cite{costa2017cnn}. These studies
are methodologically relevant because they show how a domain that is not
originally visual can become amenable to visual pattern recognition.

Source code can be treated in a similar way. A code file contains characters,
lines, indentation, punctuation, comments, delimiters, blocks, and repeated
structures. When mapped to pixels, these elements form patterns that may reflect
artifact roles. For example, a JSON file often appears as nested punctuation and
short key-value lines; an SVG file may contain dense XML-like blocks; a Java
entity may contain annotations, field declarations, and method blocks; a test
class may have repeated setup and assertion structures. Even after letters and
digits are anonymized, layout and symbolic density may remain informative.

CV4Code explicitly explored source code understanding via visual code
representations \cite{shi2023cv4code}. This line of work is closely related to
the present project because both approaches challenge the assumption that source
code must always be processed first as text, syntax trees, or graphs. The
minimap approach adds emphasis on compact rendering, anonymization, dataset
construction, explainability, and IDE integration.

\section{Texture Descriptors and Local Binary Patterns}

Before deep learning became dominant in computer vision, many image-analysis
tasks used handcrafted descriptors. Texture descriptors are particularly
relevant when an image is characterized by repeated local patterns rather than
by objects with clear semantic boundaries. Local Binary Patterns (LBP) are a
well-known texture descriptor that compares each pixel with its neighborhood and
encodes local intensity relations as binary patterns
\cite{maenpaa2005texture,martins2013database}.

In a minimap, the discriminative signal may not be an object such as a face,
vehicle, or organ. It may instead be a pattern of indentation, delimiter density,
line length, repeated symbols, or blank-space distribution. This makes LBP a
reasonable baseline for initial experiments. The first study in this doctoral
project followed this direction by extracting LBP histograms from source code
minimaps and using classical classifiers such as K-Nearest Neighbors, Random
Forest, and Support Vector Machines \cite{gebara2025iwssip}.

Handcrafted descriptors have advantages and limitations. They are simpler,
faster, and more interpretable than large neural models in many contexts.
However, they may fail to capture multi-scale or task-specific patterns that
deep models can learn directly from data. This motivates the transition from LBP
experiments to convolutional and transfer-learning experiments in the larger
study.

\section{Deep Learning and Transfer Learning}

Deep learning models can learn hierarchical representations from images. In
computer vision, convolutional neural networks such as ResNet
\cite{he2016resnet} and EfficientNet \cite{tan2019efficientnet} have become
standard backbones for image classification and feature extraction. Vision
Transformers and hybrid architectures further expanded the range of possible
models by adapting attention mechanisms to image patches \cite{vaswani2017attention}.

Transfer learning is particularly useful when a model pretrained on a large
dataset can be adapted to a smaller or more specialized dataset. Although
natural-image pretraining is not designed for source code minimaps, pretrained
filters may still provide useful low-level visual features such as edges,
textures, contrast patterns, and spatial hierarchies. The large-scale minimap
study evaluates both from-scratch and pretrained variants, showing that
pretraining can provide consistent but sometimes modest gains
\cite{gebara2025ciarp}.

The use of deep learning in this project introduces methodological concerns.
First, source code datasets are often imbalanced: some file types, projects, or
authors dominate the corpus. Second, model performance must be interpreted with
metrics beyond accuracy, especially when minority classes matter. Third,
training and evaluation splits must avoid leakage, for example when files from
the same project or author appear across partitions in ways that inflate
performance. These concerns are addressed later in the methodology and research
plan chapters.

\section{Class Imbalance in Software Datasets}

Software datasets are rarely balanced. Some file types occur frequently because
they are part of the dominant language of a repository. Some projects are much
larger than others. Some contributors or automated accounts touch many files,
while others contribute only a small number of artifacts. This long-tail
structure appears naturally in software repositories and directly affects model
evaluation.

Class imbalance can make a model appear strong while hiding weak behavior on
minority classes. For example, if a dataset is dominated by Java, C, or
JavaScript files, a model can achieve high accuracy by learning the largest
classes and ignoring rare but important artifact types. This is particularly
problematic in software engineering because rare classes may be precisely the
ones that matter in maintenance: unusual configuration files, scripts, generated
artifacts, migration files, security-sensitive files, or framework-specific
components.

The literature on imbalanced data recommends paying attention to both training
strategies and evaluation metrics \cite{he2009imbalance,japkowicz2002}. Metrics
such as macro-F1 treat all classes equally and therefore expose poor performance
on minority classes. Weighted-F1 reflects the observed distribution and is useful
for understanding practical behavior under the real dataset. ROC-AUC can help
assess separability, but it should not be used alone. Ortigosa-Hernandez et al.
\cite{ortigosa2017} also emphasize that the extent of multi-class imbalance
should be measured rather than assumed.

In minimap-based analysis, imbalance has an additional interpretive dimension.
Head classes may contain more visual diversity simply because they have more
examples. Tail classes may look visually coherent but still be difficult to
learn because the model sees too few instances. Conversely, a small class may be
easy if it has a distinctive template. For this reason, the thesis must report
both aggregate metrics and per-class behavior. Confusion matrices, class-wise
recall, and attribution maps are necessary to understand whether the model is
learning useful patterns or mostly reproducing dataset frequency.

\section{IDE Workflows and Developer Feedback}

Integrated development environments are not passive text editors. They combine
navigation, syntax highlighting, build integration, refactoring, testing,
debugging, static analysis, version-control support, and extension mechanisms.
Any tool that aims to support developers must fit into this environment without
interrupting the flow of work.

This matters for minimap-based analysis because a visual signature is useful
only if it can be generated and interpreted at the right moment. During code
review, a developer may want to inspect whether a new file resembles existing
project conventions. During maintenance, a developer may want to identify
unusual artifacts in a folder. During onboarding, a developer may want to obtain
a quick overview of file roles in an unfamiliar repository. These scenarios
require different interaction designs. A single-file command may be enough for
quick inspection, while project-level analysis may require dashboards, grouping,
filtering, and explanation overlays.

Developer feedback should also be calibrated. A tool should avoid presenting a
prediction as an unquestionable fact. Instead, it should communicate that a file
visually resembles a class, that the confidence is limited, or that the model
focused on particular regions. This is where explainability and interface design
meet. Attribution maps, confidence scores, and textual summaries can help
developers decide whether to trust, ignore, or investigate a prediction.

The MinimapSignatures prototype is therefore not only an implementation detail.
It is part of the research strategy. It allows the project to study how minimap
analysis can be delivered in IDEs and which forms of feedback are plausible for
future evaluation.

\section{Explainable Artificial Intelligence}

Explainable artificial intelligence (XAI) studies methods for making model
behavior more transparent to humans. XAI is important when models are used in
high-stakes or interpretive contexts, but it is also useful in research settings:
it helps verify whether a model is learning plausible evidence or exploiting
spurious artifacts. Surveys such as \citeonline{adadi2018xai} discuss the
broader XAI landscape, including model-specific and model-agnostic methods.

In this project, XAI is used to interpret convolutional models trained on
minimap images in the large-scale study. Integrated Gradients is especially
relevant because it attributes prediction scores to input pixels by integrating
gradients along a path from a baseline (black image) to the actual minimap input
\cite{integratedgradients2017}. When applied to minimaps, the resulting
attribution maps highlight whether a model is relying on indentation bands,
delimiter regions, dense blocks, headers, boilerplate, or other structural
regularities.

These maps are interpretive aids, not proofs of causality. A high-attribution
region indicates that the model is sensitive to that area, but it does not
necessarily mean that region carries the ground-truth discriminative signal for
the task. XAI helps connect quantitative results to qualitative understanding
and can expose risks: if Type classification relies mainly on file headers or
boilerplate rather than structural body patterns, predictions may generalize
poorly to repositories with different templates.

\section{Anonymization as Intellectual Property Protection}

A key barrier to applying external analysis tools on industrial codebases is
that source code often encodes \textit{strategic business information}: formulas
embedded in arithmetic expressions, magic constants that encode pricing or
margin logic, variable and function names that reveal domain concepts, and
comments that describe business rules. Consider a hypothetical expression such
as \texttt{(Math.sin(x) + Math.cos(y)) * 23.45}: even without knowing the
context, a reader who sees the identifier \texttt{profitMargin} and the constant
\texttt{23.45} together can infer something about the company's internal
calculations. Sending such code to an external service---even for legitimate
analysis---exposes intellectual property that the organization may not wish to
share.

The character-level anonymization strategy adopted in this project addresses
this concern directly. By mapping all letters to one of two canonical values
(uppercase or lowercase) and all digits to a single canonical value, the
transformation guarantees that:

\begin{itemize}
    \item identifier names (\texttt{profitMargin}, \texttt{lucro},
    \texttt{calculateDiscount}) become visually indistinguishable from one
    another;
    \item numeric literals (\texttt{23.45}, \texttt{0.0031}, \texttt{1000000})
    are all mapped to the same intensity and carry no magnitude information;
    \item string literals and comments lose their lexical content entirely.
\end{itemize}

What \textit{is} preserved are structural signals: indentation depth, line
length, punctuation density, the relative positions of operators and delimiters,
and the block organization of the file. These signals reflect \textit{how}
code is organized, not \textit{what} it computes or \textit{why}. This
distinction makes it viable to send an anonymized minimap to an external backend
for structural analysis without disclosing the business logic embedded in the
lexical content.

\paragraph{Author prediction as structural validation.}
The Author prediction objective plays a specific methodological role in this
context. If a model can reliably identify the stylistic patterns of a
contributor even after all identifiers, literals, and comments are suppressed,
this constitutes direct evidence that the information captured by the minimap
is \textit{purely structural}: it comes from indentation habits, spacing
conventions, punctuation placement, and code organization---not from names or
values. Author prediction is therefore not presented as a deployment target but
as a \textit{proof-of-concept for the anonymization argument}: if style
survives suppression, then the model genuinely operates on structure, and
business-sensitive content is genuinely excluded.

\paragraph{Limits of anonymization.}
Anonymization in this sense is not a formal privacy guarantee against all
possible attacks. A highly distinctive architectural pattern or a very unusual
structural idiom could still be recognizable. The claim is more modest and more
useful: the approach removes the lexical layer that carries direct business
information, making the representation appropriate for contexts where the
concern is intellectual property exposure rather than individual privacy.
Stronger transformations---such as removing all punctuation or imposing uniform
indentation---can be evaluated as a trade-off between suppression strength and
predictive performance, and this is part of the planned research.

\section{Empirical Evaluation in Software Engineering}

Because this project proposes and evaluates a software-engineering method, it
must also follow principles from empirical software engineering. Guidelines for
systematic literature reviews \cite{kitchenham2007slr}, controlled and empirical
software-engineering studies \cite{wohlin2012experimentation,kitchenham2002preliminary},
case-study research \cite{runeson2012case}, and families of experiments
\cite{basili1999families} provide useful foundations for planning evaluation and
reporting threats to validity.

The empirical nature of the project appears in three places. First, the
literature mapping must be transparent about search, filtering, grouping, and
interpretation. Second, the model evaluation must describe datasets, splits,
metrics, baselines, and class imbalance. Third, the tool-support evaluation must
distinguish functional validation from usability, usefulness, or productivity
claims. A plugin that successfully sends a minimap to a backend and displays a
prediction is functionally validated; it is not yet proven to improve developer
performance. This distinction is important for the remaining doctoral plan.
